Electroencephalography (EEG) offers a non-invasive lens into human brain activity, but building large‐scale models is hampered by \emph{topological heterogeneity}: each public EEG data defines its own electrode layout, limiting generalization. We introduce \textbf{LUNA} (\textbf{L}atent \textbf{U}nified \textbf{N}etwork \textbf{A}rchitecture), a self-supervised foundation model that reconciles disparate electrode geometries while scaling linearly---not quadratically---with channel count. LUNA compresses multi-channel EEG into a fixed-size, topology-agnostic latent space via \emph{learned queries} and cross-attention. Downstream transformer blocks then operate exclusively on this latent representation using patch-wise temporal self-attention, decoupling computation from electrode count. Pre-trained on TUEG and Siena ($>$ 21,000 hours of raw EEG across diverse montages) using a masked-patch reconstruction objective, LUNA transfers effectively to four downstream tasks: abnormality detection, artifact rejection, slowing classification, and emotion recognition. It demonstrates highly competitive performance across several benchmarks, achieving state-of-the-art results on TUAR and TUSL, e.g., \textbf{0.921 AUROC} on TUAR, while reducing FLOPs by \textbf{300$\times$} and trimming GPU memory use by up to \textbf{10$\times$}. Critically,  these gains are consistent across all evaluated electrode configurations. Code is available at \url{https://github.com/pulp-bio/biofoundation}